\needspace{5\baselineskip}
\section{创新点与总结展望}

\subsection{技术创新点}

\subsubsection{按“可核验程度”分层的图谱模型}

多数文献图把引用、语义与聚类混在一张扁平表里，结果是模型幻觉与真实引用看起来一样可信。本项目按可核验程度把图谱分成三层：L1 结构层只存上游事实，L2 语义层由模型生产但每条边必须带原文证据，L3 派生层由算法生成且可复算。三层依赖单向，L2 读 L1，L3 读 L1+L2，任何一层都不得反向改写下层。这个模型不是“数据库表设计”，而是产品可信度的地基：用户看到的每一种线型、颜色与证据卡，都能在模型里找到明确位置。

\subsubsection{硬约束前置的多标签关系判定}

关系判定采用“确定性硬约束过滤 + 多标签 rubric 判定 + 引用片段核验”的两段式管道。硬约束先把候选从 $O(n^2)$ 砍到 $O(n\cdot k)$，不花 token 也不会幻觉；多标签判定让同一对论文在一次调用里比较所有可判关系，避免重复送摘要与 rubric；引用核验把“模型说”变成“原文说”，核验失败的边直接丢弃并计入 \code{rejectedQuotes}。反例与正例一样是用户可编辑的数据，这让“同期”与“思路类似”这类容易互相污染的关系有了明确边界。

\subsubsection{坐标持久的增量星空}

新增种子时，系统不重排整张图：老节点坐标作为锚点参与增量松弛，新节点先插到邻居重心附近再自由松弛；社区重聚类后按新旧社区重叠票数做编号对齐，让大部分论文保持原来的社区编号。世界原点固定在中心，世界只做对称长大。这个设计解决的是“我一搜，原来的星图就没了”的真实问题，也让“同一片星空反复并入种子”从演示噱头变成可用的工作台。

\subsubsection{模型只规划、规则判停止的 Agent 闭环}

站内 Agent 的五个工具、两轮、五步都是硬预算；模型只决定“下一步查什么”，停止由覆盖率、缺口、轮次、步数与“是否有新增证据”这些确定性规则判定。每个决策生成 \code{AskDecision}，包含停止码、理由、缺口与统计，界面可完整回放。这样既避免了“模型觉得自己答完了”的不可复现，也让成本可控：证据足够时不再多花一次调用。

\subsubsection{混合检索与图结构扩展}

检索不是“向量一刀切”：BM25 负责词面与专有名词，向量负责同义与跨语言表达，RRF 把两路名次融合，图结构扩展把“问题提到的论文之间的关系证据”补进来，最后按问题覆盖、论文命中、章节意图与图证据重排。向量不可用时自动退回 BM25，并在诊断中如实标注；这让系统在“没有 embedding 配额”的演示环境里仍然可用。

\subsubsection{语义缩放的 Canvas 渲染管线}

Canvas 2D 自写渲染管线把上千节点/边的绘制压成一次调用；布局在服务端完成，前端只做投影；悬浮/选中双写避免高频状态触发重渲染；动画按秒计算并封顶 \code{dt}。语义缩放不是“把同一张图放大缩小”，而是在领域、社区、论文三个抽象层级之间切换，阈值集中在 \code{VIEW_LEVEL_THRESHOLDS}，不散落在组件里。渲染函数是纯函数，可脱离 React 测试，也可在 Node 里生成截图用于文档。

\subsubsection{只读 MCP 与按账号隔离的密钥体系}

9 个 MCP 工具把星空开放给外部 Agent，但全部只读：改图必须由用户确认，不能由一个自己乱逛的 Agent 决定。个人令牌只在创建时返回一次，数据库只存 SHA-256，每个令牌绑定一个 \code{owner_id}；HTTP 与 stdio 共用同一份工具定义。模型 API Key 按账号绑定、AES-256-GCM 加密落库、对外只返回尾号。这套体系让“多厂商可切换”与“多账号隔离”同时成立，而不是牺牲其中一个。

\subsection{应用创新点}

\begin{enumerate}
  \item \textbf{把文献调研从“读列表”变成“看星空”。} 论文是星，关系是引力线；拉远看领域星云，推近看论文星体。空间记忆比列表记忆更符合人对“方向感”的需求。
  \item \textbf{让关系可以被质疑。} 每条语义边都有判定标准、反例、置信度与原文证据；用户不需要相信一个分数，只需要读一眼原文。这是与“直接给结论”的问答产品最根本的区别。
  \item \textbf{让调研可以累积。} 同一片星空可以反复并入新种子，老星体不漂移；历史星空保存视角与缩略预览，用户可以回到上次的研究现场继续追问。
  \item \textbf{让图表可以提问。} 上传截图、表格或曲线，系统抽取 OCR、表格行列、图表数值与区域坐标，回答里的角标可以点回图中的具体位置。
  \item \textbf{让研究现场可以被外部 Agent 复用。} MCP 把星空、论文、关系证据与检索能力开放成标准工具；外部 Agent 不必重新抓取、重新判定，就能在用户已保存的研究上下文里工作。
  \item \textbf{降低跨学科调研的门槛。} 对不熟悉某个方向术语的读者，“基线 / 同期 / 同任务异方法 / 结论矛盾”比一堆引用数更接近他们真正想问的问题；这具有开放科学与科研公平的公共价值。
\end{enumerate}

\subsection{总结}

StellarLit 完成了一条从“种子论文”到“可追问研究星图”的完整链路：种子解析永远让用户确认；建图链路把引文邻域、结构事实、语义关系、社区与布局组织成三层图谱；问答链路用混合检索、Agent 补查与证据覆盖判定给出带角标的回答；多模态把图表变成可定位的证据；MCP 把星空开放给外部 Agent。工程上，Next.js 15 + TypeScript strict + PostgreSQL + Redis + Docker Compose 的组合让这套能力可以在一台 2 核 4G 的服务器上跑起来，15 项自检让“能演示”变成“能重复”。

更重要的不是功能数量，而是取舍：我们没有让模型自由发挥，而是把“边界”交给确定性规则，把“语言理解”交给模型，把“可信”交给证据核验，把“成本”交给缓存、批处理与角色拆分。这个取舍贯穿图谱模型、关系判定、RAG、Agent 与安全设计，也是本项目最想表达的技术观点。

\subsection{局限与展望}

\begin{table}[htbp]
  \centering
  \small
  \caption{当前局限与后续路线}
  \label{tab:future}
  \begin{tabularx}{\textwidth}{L{2.8cm}L{5.6cm}Y}
    \toprule
\textbf{局限} & \textbf{当前状态} & \textbf{下一步} \\
    \midrule
    证据粒度 & 语义证据到摘要与关系原文，尚未按章节切分全文 & 接入 arXiv 全文，按章节与字符偏移切块；引用精确到“第几节第几段” \\
    评测体系 & 以自检与探针为主，缺少公开标注集上的定量指标 & 构建领域标注集，报告关系 precision/recall、RAG 命中率与证据忠实度 \\
    多模态范围 & 覆盖文本与图像，语音与视频尚未接入 & 浏览器语音输入/朗读；视频抽帧后作为图像证据 \\
    枢纽种子的“同期” & 高被引论文的邻域几乎都引用了它，按定义“同期”通常为空 & 在 UI 中解释这是定义结果；提供“放宽互不引用”的可选视图 \\
    多实例限流 & IP/账号窗口计数在单进程内存，重启归零 & 迁移到 Redis 计数；多实例共享同一限流桶 \\
    协作能力 & 单账号工作台，无团队共享与评论 & 增加星空共享、批注与版本对比；权限模型复用现有 owner 隔离 \\
    \bottomrule
  \end{tabularx}
\end{table}

路线图上最优先的是“全文证据”与“评测体系”。前者让引用从摘要级走到章节级，后者让“关系判定准不准、RAG 答得对不对”从主观感受变成可比较的数字。再往后是语音/视频、多实例部署与团队协作。我们希望 StellarLit 最终不只是一个比赛作品，而是一个研究者愿意长期打开的工作台：星空可以慢慢长大，问题可以一层层追问，而每一条关系都经得起点开看一眼。
